\documentclass{article}
\usepackage[utf8]{inputenc}
\usepackage{amsmath}
\usepackage{geometry}
\geometry{margin=2cm}
\begin{document}

\section*{Análise da Questão 158 – ENEM 2024 – Matemática}

\subsection*{Enunciado}
Em uma loja de defensivos agrícolas, alguns preços foram divulgados em um cartaz:
\begin{itemize}
  \item Defensivo Tipo A: R$ 4,20 por litro, com promoção para compra de ao menos 35 litros (máscara grátis);
  \item Defensivo Tipo B: R$ 3,00 por litro, máscara custa R$ 12,50 (sem promoção).
\end{itemize}
Sabe-se que:
\begin{itemize}
  \item 1 litro do Tipo A é suficiente para 0,5 hectare;
  \item 1 litro do Tipo B é suficiente para 0,4 hectare.
\end{itemize}
Um agricultor deseja adquirir defensivo suficiente para aplicação em 20 hectares, além de levar uma máscara para aplicação.

\subsection*{Solução Técnica}

Para determinar o menor custo de compra, calculamos a quantidade de litros necessária para cada tipo, os custos totais incluindo a máscara (quando aplicável) e comparamos.

\paragraph{Quantidades necessárias:}
\begin{align*}
\text{Litros Tipo A} &= \frac{20}{0{,}5} = 40 \text{ litros}, \\
\text{Litros Tipo B} &= \frac{20}{0{,}4} = 50 \text{ litros}.
\end{align*}

\paragraph{Custos totais:}
\begin{itemize}
  \item Tipo A: 40 litros, promove máscara grátis (compra \geq 35 litros).
  \begin{align*}
  C_A &= 40 \times 4{,}20 = R\$ 168{,}00.
  \end{align*}

  \item Tipo B: 50 litros, máscara paga separadamente.
  \begin{align*}
  C_B &= 50 \times 3{,}00 + 12{,}50 = R\$ 162{,}50.
  \end{align*}

  \item Misturas não compensam devido a custo extra da máscara se Tipo A for comprado em menos de 35 litros.
\end{itemize}

\paragraph{Conclusão:} O menor valor gasto é R$ 162,50, adquirindo 50 litros do Tipo B e a máscara.

\subsection*{Explicação Didática}

\begin{enumerate}
  \item Calcule a quantidade de defensivo necessária para 20 hectares com cada tipo.
  \item Verifique se a compra do Tipo A alcança a promoção da máscara grátis.
  \item Calcule o custo total para cada cenário.
  \item Compare os custos e escolha o mínimo.
  \item Reforce que misturar produtos com custo de máscara pode não reduzir o gasto.
\end{enumerate}

\subsection*{Distratores}
\begin{itemize}
  \item R$ 147,00: subestimação da quantidade necessária de Tipo A.
  \item R$ 150,00: esquecer custo da máscara no Tipo B.
  \item R$ 165,75: erro de cálculo ou mistura mal avaliada.
  \item R$ 168,00: custo correto para 40 litros do Tipo A, mas não o mínimo custo.
\end{itemize}

\subsection*{Perfil da Questão}

Trata-se de uma questão de nível médio, que exige habilidades de raciocínio matemático, interpretação de texto, cálculo de proporção e análise de promoções comerciais, com interpretação visual do cartaz.

\newpage

\subsection*{Referência Visual (Descrição do Cartaz)}
\begin{itemize}
  \item Promoção: compra de ao menos 35 litros do Tipo A dá direito à máscara para aplicação.
  \item Preços: Tipo A custa R$ 4,20 por litro; Tipo B custa R$ 3,00 por litro; máscara custa R$ 12,50 (se comprada separadamente).
\end{itemize}

Este conjunto sistematiza os elementos do problema e a solução, clarificando o raciocínio e facilitando o entendimento didático da questão.

\end{document}